\documentclass[10pt,a4paper,ragged2e,withhyper,normalphoto]{files/altacv}
% Change the page layout if you need to
\geometry{left=1.25cm,right=1.25cm,top=1.5cm,bottom=1.5cm,columnsep=1.25cm}
\usepackage{paracol}
\usepackage{marvosym}
\usepackage{booktabs}
\usepackage[backend=biber,style=ieee,sorting=ydnt,defernumbers=true]{biblatex}
% Font selection for pdflatex
\usepackage[rm]{roboto}
\usepackage[defaultsans]{lato}
% Start Hyphenation Rules
% -----------------------------------------------------------------------------
% Line Breaking & Hyphenation
%
% These settings reduce automatic word hyphenation while keeping text within
% the page margins.
%
% \hyphenpenalty=10000
%   Prevents TeX from automatically splitting words (e.g., "organi-zational").
%   0 → Hyphenation is allowed freely.
%   50 → Default (approximately).
%   1000 → Avoid hyphenation when possible.
%   10000 → Never hyphenate automatically.
%
% \exhyphenpenalty=10000
%   Prevents line breaks at explicit hyphens in words such as
%   "state-of-the-art".
%
% \emergencystretch=3em
%   Gives TeX additional flexibility to stretch inter-word spaces when
%   necessary. This helps move an entire word to the next line instead of
%   hyphenating it, while reducing overfull lines.
%
% Higher values reduce hyphenation but may produce slightly wider spaces
% between words in difficult paragraphs.
% -----------------------------------------------------------------------------
\hyphenpenalty=10000
\exhyphenpenalty=10000
\emergencystretch=1em
% End of Hyphenation Rules
\renewcommand{\familydefault}{\sfdefault}

% Indigo theme colors
\definecolor{primary}{RGB}{63,81,181}
\definecolor{secondary}{RGB}{50,50,50}
\colorlet{name}{black}
\colorlet{tagline}{primary}
\colorlet{heading}{primary}
\colorlet{headingrule}{primary}
\colorlet{subheading}{secondary}
\colorlet{accent}{primary}
\colorlet{emphasis}{secondary}
\colorlet{body}{secondary}
\definecolor{subtle}{RGB}{100,100,100}

% Font and formatting settings
\renewcommand{\namefont}{\Large\rmfamily\bfseries}
\renewcommand{\personalinfofont}{\footnotesize}
\renewcommand{\cvsectionfont}{\large\rmfamily\bfseries}
\renewcommand{\cvsubsectionfont}{\large\bfseries}

% Section formatting
% -----------------------------------------------------------------------------
% It reduces the height of the underline
% of the section title
%
% Default: 2pt
% -----------------------------------------------------------------------------
\let\oldcvsection\cvsection
\renewcommand{\cvsection}[2][]{%
\nointerlineskip\bigskip%
\ifstrequal{#1}{}{}{\marginpar{\vspace*{\dimexpr1pt-\baselineskip}\raggedright\input{#1}}}%
{\color{heading}\cvsectionfont\MakeUppercase{#2}}\\[-1ex]%
{\color{headingrule}\rule{\linewidth}{0.5mm}\par}\medskip
}
% End of section formatting

\begin{document}

\name{Md. Fahad Hossain}
\tagline{}
\photoR{2.4cm}{files/Md_Fahad_Hossain.jpeg}
\personalinfo{
    \phone{+8801732-327112} \\[1mm]
    \email{hossainfahad725@gmail.com} \\[1mm]
    \location{South Betduba, Kalihati, Kalihati, Tangail}
}
\makecvheader
% \divider


\cvsection{Career Objective}
To build a career enriched with the capability to adapt different organizational cultures of environments through the knowledge. I have gained academic qualifications with my honesty, punctuality and effective communication skills as well.
\vspace{2mm}
\par


\cvsection{Experience}

{\large\color{emphasis}Sales Officer\par}
\smallskip\normalsize
\textbf{\color{accent}Unilever Bangladesh Limited}\hfill{\small\color{body}March 2019 -- April 2023}\par

\divider

{\large\color{emphasis}Relationship Officer\par}
\smallskip\normalsize
\textbf{\color{accent}Youth for Banglaesh}\hfill{\small\color{body}November 2024 -- Present}\par
\vspace{2mm}
\par

\cvsection{Education}

\setlength{\tabcolsep}{8pt}
\renewcommand{\arraystretch}{1.5}
\begin{tabularx}{\linewidth}{|>{\raggedright\arraybackslash}p{4.2cm}|>{\raggedright\arraybackslash}X|>{\raggedright\arraybackslash}p{2.2cm}|r|}
    \hline
    \textbf{Degree} & \textbf{Institution} & \textbf{Year} & \textbf{Result} \\
    \hline
    Bangla Master's Degree & GOVT. SAADAT COLLEGE & 2015 -- 2021 & CGPA: 2.75/4.00 \\
    \hline
    Bangla Honour's Degree & GOVT. SAADAT COLLEGE & 2015 -- 2020 & CGPA: 2.78/4.00 \\
    \hline
    Higher Secondary School Certificate (H.S.C) & HAZI ABU HASHEM TECHNICAL \& BUSINESS MANAGEMENT COLLEGE & 2016 & GPA: 4.50/5 (Group: B.M.) \\
    \hline
    Secondary School Certificate (S.S.C) & GOVT. TEXTAIL INSTITUTE, KALIHATI, TANGAIL & 2014 & GPA: 4.85/5 \\
    \hline
\end{tabularx}
\vspace{2mm}
\par

\cvsection{Skills}
\newcommand{\skilltag}[2]{\textbf{#1:} #2}

\skilltag{Microsoft Office}{MS Word, MS Excel \& MS PowerPoint}\\[0.5em]
\skilltag{Others}{Email \& Internet Browsing}

\vspace{2mm}

\cvsection{Personal Details}
\renewcommand{\arraystretch}{1.6}
\begin{tabular}{@{}>{\raggedright\arraybackslash}p{4.5cm} l@{}}
    \textbf{Father's Name}  & : Md Anowar Hossain                \\
    \textbf{Mother's Name}  & : Fahima Khan                      \\
    \textbf{Date Of Birth}  & : 06-05-1998                       \\
    \textbf{Place Of Birth} & : Tangail                          \\
    \textbf{Sex}            & : Yes                              \\
    \textbf{Marital Status} & : Married                          \\
    \textbf{Nationality}    & : Bangladeshi                      \\
    \textbf{Religion}       & : Islam                            \\
    \textbf{Blood Group}    & : A+                               \\
    \textbf{Height}         & : 5\textquotedbl 8\textquotesingle \\
\end{tabular}
\vspace{2mm}
\par


\end{document}
